\documentclass[11pt]{article}

% Change "review" to "final" to generate the final (sometimes called camera-ready) version.
% Change to "preprint" to generate a non-anonymous version with page numbers.
\usepackage[final]{acl}

% Standard package includes
\usepackage{times}
\usepackage{latexsym}

% For proper rendering and hyphenation of words containing Latin characters (including in bib files)
\usepackage[T1]{fontenc}
% For Vietnamese characters
% \usepackage[T5]{fontenc}
% See https://www.latex-project.org/help/documentation/encguide.pdf for other character sets

% This assumes your files are encoded as UTF8
\usepackage[utf8]{inputenc}

% This is not strictly necessary, and may be commented out,
% but it will improve the layout of the manuscript,
% and will typically save some space.
\usepackage{microtype}

% This is also not strictly necessary, and may be commented out.
% However, it will improve the aesthetics of text in
% the typewriter font.
\usepackage{inconsolata}

%Including images in your LaTeX document requires adding
%additional package(s)
\usepackage{graphicx}

% If the title and author information does not fit in the area allocated, uncomment the following
%
%\setlength\titlebox{<dim>}
%
% and set <dim> to something 5cm or larger.

\title{What is the price of a model's eye?:\\ Understanding How VLMs Answer Multiple Choice Questions in Multimodal Settings}

% Author information can be set in various styles:
% For several authors from the same institution:
% \author{Author 1 \and ... \and Author n \\
%         Address line \\ ... \\ Address line}
% if the names do not fit well on one line use
%         Author 1 \\ {\bf Author 2} \\ ... \\ {\bf Author n} \\
% For authors from different institutions:
% \author{Author 1 \\ Address line \\  ... \\ Address line
%         \And  ... \And
%         Author n \\ Address line \\ ... \\ Address line}
% To start a separate ``row'' of authors use \AND, as in
% \author{Author 1 \\ Address line \\  ... \\ Address line
%         \AND
%         Author 2 \\ Address line \\ ... \\ Address line \And
%         Author 3 \\ Address line \\ ... \\ Address line}

% For several authors from the same institution:
\author{Divyanshu Giri, Rachit Mehta, Yashav Bhatnagar \\
        IIIT, Hyderabad}

% \author{Divyanshu Giri \\
%   Affiliation / Address line 1 \\
%   Affiliation / Address line 2 \\
%   Affiliation / Address line 3 \\
%   \texttt{divyanshu.giri@research.iiit.ac.in} \\\And
%   Rachit Mehta \\
%   Affiliation / Address line 1 \\
%   Affiliation / Address line 2 \\
%   Affiliation / Address line 3 \\
%   \texttt{rachit.mehta@students.iiit.ac.in} \\}

%\author{
%  \textbf{First Author\textsuperscript{1}},
%  \textbf{Second Author\textsuperscript{1,2}},
%  \textbf{Third T. Author\textsuperscript{1}},
%  \textbf{Fourth Author\textsuperscript{1}},
%\\
%  \textbf{Fifth Author\textsuperscript{1,2}},
%  \textbf{Sixth Author\textsuperscript{1}},
%  \textbf{Seventh Author\textsuperscript{1}},
%  \textbf{Eighth Author \textsuperscript{1,2,3,4}},
%\\
%  \textbf{Ninth Author\textsuperscript{1}},
%  \textbf{Tenth Author\textsuperscript{1}},
%  \textbf{Eleventh E. Author\textsuperscript{1,2,3,4,5}},
%  \textbf{Twelfth Author\textsuperscript{1}},
%\\
%  \textbf{Thirteenth Author\textsuperscript{3}},
%  \textbf{Fourteenth F. Author\textsuperscript{2,4}},
%  \textbf{Fifteenth Author\textsuperscript{1}},
%  \textbf{Sixteenth Author\textsuperscript{1}},
%\\
%  \textbf{Seventeenth S. Author\textsuperscript{4,5}},
%  \textbf{Eighteenth Author\textsuperscript{3,4}},
%  \textbf{Nineteenth N. Author\textsuperscript{2,5}},
%  \textbf{Twentieth Author\textsuperscript{1}}
%\\
%\\
%  \textsuperscript{1}Affiliation 1,
%  \textsuperscript{2}Affiliation 2,
%  \textsuperscript{3}Affiliation 3,
%  \textsuperscript{4}Affiliation 4,
%  \textsuperscript{5}Affiliation 5
%\\
%  \small{
%    \textbf{Correspondence:} \href{mailto:email@domain}{email@domain}
%  }
%}

\begin{document}
\maketitle
\begin{abstract}
MCQA is a key competence of vision-language models that is tested by a variety of mainstream benchmarks. However, these models can have a wide range of changes in performance when the task is diversified even slightly. The question we are asking is: \textit{How do mainstream vision-language models perform on formatted multiple-choice question answering tasks?} We set out to evaluate and understand the models' internal representations by employing methods like layer-wise activation patching, individual attention head localization, and vocabulary projection for answer choices to localize key hidden states and activation patterns that encode information about predictive tokens. We will also employ novel methods like visual-answer choices, comparison of symbol binding vs. visual grounding, spatial attention analysis, as well as cross-image activation patching so as to comprehensively probe the attention and activation information encodings within a model.
\end{abstract}

\section{Introduction}

Multi-choice question answering (MCQA) has been a popular task for evaluating LMs \cite{wiegreffe2025answerassembleaceunderstanding}, and has been gaining traction in the space of vision-language models recently \cite{meng2024mmiumultimodalmultiimageunderstanding}. A prominent example is the MMMU benchmark \cite{yue2024mmmumassivemultidisciplinemultimodal} which evaluates multimodal models on massive multi-discipline tasks on 28 open-source LLMs as well as proprietary models such as GPT-4V (Vision) and Gemini. While LMs have been evaluated comprehensively in this domain, VLMs are yet to be probed in depth.

Our methodology is rooted in NLP interpretability. Once an image is encoded and projected into the language model’s embedding space, it becomes a sequence of tokens that is treated exactly like text tokens. This allows us to probe the model’s language component and treat the vision encoder as a black box. Vocabulary projection and activation patching are defined in terms of a model’s vocabulary and residual stream. Our vision specific experiments, spatial attention analysis, cross image activation patching and image as option binding, described later, all operate on image tokens in the language model’s residual stream, no intervention on the vision encoder itself happens. 

Our project extends \cite{wiegreffe2025answerassembleaceunderstanding}, an NLP interpretability paper on text-only language models. Our focus is not on whether the model answers correctly but on which layer and which attention heads produce that answer. While benchmarks like MMLU and MMMU are widely trusted, prior work shows that text-only models can be sensitive to superficial changes, such as answer order or label symbols \cite{pezeshkpour2024large,alzahrani2024when,khatun2024study,zheng2024large}. Our works falls within the line of this prior interpretability work and by adding vision as a new variable we ask whether this underlying brittleness extends to multimodal benchmarks. 

\section{Motivation}

The use of vision language models is increasing in real world environments with their capabilities being judged predominantly from benchmarks featuring formatted multi-choice question answering (MCQA) such as MMMU. While these benchmark accuracy metrics provide a broad view of model’s abilities, they treat model as a black box and fail to reveal how the model arrives at a specific answer. 

This present a significant reliability risk as recent mechanistic interpretability work on text-only language models demonstrates that models often rely on superficial heuristics such as positional bias, answer order or label symbols, rather than semantic reasoning, while getting high benchmark scores. Modern VLMs are built by using these language models as a backbone hence the question of whether VLMs ground their answers in visual evidence or do they utilise text-biased heuristics of their backbone models becomes important. A model could correctly identify an object in an image but map it to the wrong symbol answer due to inherited text processing bias and record a failure in the benchmark, or conversely, a model might guess correctly based on text symbols bias without utilizing visual input.

To ensure VLMs can be trusted in high-stakes tasks, understanding precise layers and attention mechanisms that govern multimodal MCQA becomes essential. By extending existing NLP interpretability techniques into multimodal space, we aim to understand whether model’s answers are driven by visual representations or its pre-existing textual heuristics

\section{Contributions}

Recent work by Wiegreffe et al. \cite{wiegreffe2025answerassembleaceunderstanding} studies the internal mechanisms underlying formatted MCQA in text-only language models. Using techniques such as vocabulary projection and activation patching, they identify layers and attention mechanisms involved in promoting specific answer choices, and further localize sparse attention heads responsible for answer-symbol prediction. We modify and extend this line of analysis to vision-language models, asking whether these mechanisms persist when answer selection additionally depends on visual information, and how visual information is incorporated into the answer-selection process.

Concretely, our contributions are: (i) testing whether the layer-wise and attention-head mechanisms Wiegreffe et al. identify for text-only MCQA persist across four VLM families; (ii) new diagnostics -- answer-content vs. answer-symbol localization, and position vs. content disentanglement -- for separating what a VLM represents internally from how that representation is expressed as a choice symbol; and (iii), time permitting, vision-specific extensions with no text-only analog, including image-as-option symbol binding and cross-image activation patching, which isolate the exact layer at which visual evidence enters the answer-selection pathway.

\section{Methodology}

In this project, we set out to evaluate and understand how models promote the prediction of specific answer choices, in the space of attention, MLP, and different layers themselves. We plan to evaluate the models from families such as Qwen \cite{bai2025qwen3vltechnicalreport}, PaliGemma \cite{beyer2024paligemmaversatile3bvlm}, LLaVA-OneVision \cite{li2024llavaonevisioneasyvisualtask}, and InternVL \cite{wang2025internvl35advancingopensourcemultimodal}. Our proposed methods can be summarized as follows:

\subsection{Core Methods}

\begin{enumerate}

\item \textbf{Layer-wise Activation Patching}: Patch hidden states at the final token position from one prompt format (e.g., answer at position A) into another (answer at position B), layer by layer. Identifies which layers causally encode the answer choice in VLMs, and whether the critical layers are the same as in text-only LMs.

\item \textbf{Attention vs. MLP Decomposition}: Project MHSA and MLP outputs separately into vocabulary space at each layer. Determines whether attention heads still dominate discriminative answer symbol production over MLPs in VLMs, or if the MLP plays a bigger role due to visual feature integration.

\item \textbf{Individual Attention Head Localization}: Tests whether VLM answer production is still driven by a sparse set of attention heads or requires a broader set.

\item \textbf{Vocabulary/Logit Projection for Answer Choices}: Apply logit lens layer-by-layer: project each layer's hidden state to the vocabulary and track how the logits for A/B/C/D evolve across depth. Shows the trajectory of answer symbol emergence in VLMs and whether it mirrors the smooth late-layer rise seen in text-only models.

\item \textbf{Answer-Content vs. Answer-Symbol Localization}: Separately track when the model's hidden states encode the semantic content of the correct answer (e.g., "yellow") vs. the symbol (e.g., "A"). Reveals whether VLMs separate the "what is the answer" and "what letter is it" computations into distinct layers like text LMs do.

\item \textbf{Position Information vs. Answer-Content Information}: Patch between prompts where the correct answer content stays the same but its position changes ($A \rightarrow C$). Disentangles which layers encode where the correct answer sits in the list vs. what the correct answer actually is, and whether visual grounding adds an extra processing stage.

\item \textbf{Arbitrary Answer Symbols}: Replace A/B/C/D with out-of-distribution symbols (Q/Z/R/X, 1/2/3/4) and track logit trajectories. Tests whether VLMs exhibit the same two-stage process (first promoting A/B/C/D, then switching to OOD symbols at late layers) observed in text-only models.

\end{enumerate}

\subsection{Stretch Goals}

We also intend to do these novel vision-specific experiments if the duration of the course project permits, otherwise we'll continue evaluation on our own time expense:

\begin{enumerate}

\item \textbf{Visual Answer Choices (Image-as-Option)}: Present answer choices as images instead of text (e.g., colored patches, diagrams). The model must bind a visual option to a symbol letter (no string matching possible). Tests whether the symbol binding circuit works at all when answer content is non-textual, and whether it requires more layers or different heads.

\item \textbf{Cross-Image Activation Patching}: Keep text completely fixed, but patch in image token activations from a different image (where a different answer is correct). Track at which layer the answer symbol flips. This has no text-only analog; it cleanly isolates the exact layer where visual evidence enters the symbol selection pathway, since any change in the predicted symbol is purely caused by the patched image representations.

\item \textbf{Spatial Attention Analysis (Question-Driven vs. Option-Driven)}: For questions that reference specific image regions ("What is in the top-left?"), track which image token positions the answer-critical attention heads attend to. Does the model attend to the region the question asks about, or does it attend to regions corresponding to each answer option and compare? Reveals whether visual reasoning is question-guided or option-guided.

\item \textbf{Failure Mode Decomposition (Symbol Binding vs. Visual Grounding)}: Compare layer-wise logit projection trajectories for correct predictions vs. wrong predictions. For a wrong prediction, determine whether the model first represents the wrong answer content and then correctly maps it to its corresponding option symbol $\rightarrow$ visual grounding failed but symbol binding is intact. If the correct answer content is represented but the model maps it to the wrong option symbol $\rightarrow$ symbol binding failed despite correct visual grounding. Separates upstream visual understanding failures from downstream answer-symbol binding failures.

\end{enumerate}

\section{Dataset}

We evaluate on the MMMU benchmark \cite{yue2024mmmumassivemultidisciplinemultimodal}, which spans massive multi-discipline tasks across text and visual reasoning and has been used to evaluate both open-source and proprietary vision-language models. MMMU items provide the base pool from which we construct, for each question, the four position variants and the alternate-symbol variants (\texttt{Q/Z/R/X}, \texttt{1/2/3/4}) used in the per-model-family screening procedure and in the layer-wise activation patching experiments described above.

\section{Evaluation Plan}

Before applying any interpretability method, we first screen each model family for the model that answers formatted MCQA \textit{consistently}, following \cite{wiegreffe2025answerassembleaceunderstanding}: for each item we construct four versions with the correct answer at each of the four positions, plus versions using unusual answer symbols (\texttt{Q/Z/R/X}, \texttt{1/2/3/4}), and take the minimum accuracy across all of these variants. This penalizes a model whose apparent accuracy comes from a positional or symbol-frequency bias rather than genuine answer selection, and we keep only the strongest model per family, requiring at least 20\% above random accuracy, before proceeding. For activation patching, we only compute effects on instances a model predicts correctly under both prompt versions being compared, so that causal effects are not confounded with cases the model never understood in the first place. At each layer, we report two complementary metrics: logit values (the score assigned to each answer symbol, the gap between the top and second-choice symbol, and the largest score assigned to any other vocabulary token as a baseline) and probits (softmax-normalized probabilities). Logits can reveal a component actively suppressing an answer choice, an effect that can vanish after normalization, while probits let us compare confidence fairly across our four model families, which have different raw logit scales.

% \section{Engines}

% To produce a PDF file, pdf\LaTeX{} is strongly recommended (over original \LaTeX{} plus dvips+ps2pdf or dvipdf).
% The style file \texttt{acl.sty} can also be used with
% lua\LaTeX{} and
% Xe\LaTeX{}, which are especially suitable for text in non-Latin scripts.
% The file \texttt{acl\_lualatex.tex} in this repository provides
% an example of how to use \texttt{acl.sty} with either
% lua\LaTeX{} or
% Xe\LaTeX{}.

% \section{Preamble}

% The first line of the file must be
% \begin{quote}
% \begin{verbatim}
% \documentclass[11pt]{article}
% \end{verbatim}
% \end{quote}

% To load the style file in the review version:
% \begin{quote}
% \begin{verbatim}
% \usepackage[review]{acl}
% \end{verbatim}
% \end{quote}
% For the final version, omit the \verb|review| option:
% \begin{quote}
% \begin{verbatim}
% \usepackage{acl}
% \end{verbatim}
% \end{quote}

% To use Times Roman, put the following in the preamble:
% \begin{quote}
% \begin{verbatim}
% \usepackage{times}
% \end{verbatim}
% \end{quote}
% (Alternatives like txfonts or newtx are also acceptable.)

% Please see the \LaTeX{} source of this document for comments on other packages that may be useful.

% Set the title and author using \verb|\title| and \verb|\author|. Within the author list, format multiple authors using \verb|\and| and \verb|\And| and \verb|\AND|; please see the \LaTeX{} source for examples.

% By default, the box containing the title and author names is set to the minimum of 5 cm. If you need more space, include the following in the preamble:
% \begin{quote}
% \begin{verbatim}
% \setlength\titlebox{<dim>}
% \end{verbatim}
% \end{quote}
% where \verb|<dim>| is replaced with a length. Do not set this length smaller than 5 cm.

% \section{Document Body}

% \subsection{Footnotes}

% Footnotes are inserted with the \verb|\footnote| command.\footnote{This is a footnote.}

% \subsection{Tables and figures}

% See Table~\ref{tab:accents} for an example of a table and its caption.
% \textbf{Do not override the default caption sizes.}

% \begin{table}
%   \centering
%   \begin{tabular}{lc}
%     \hline
%     \textbf{Command} & \textbf{Output} \\
%     \hline
%     \verb|{\"a}|     & {\"a}           \\
%     \verb|{\^e}|     & {\^e}           \\
%     \verb|{\`i}|     & {\`i}           \\
%     \verb|{\.I}|     & {\.I}           \\
%     \verb|{\o}|      & {\o}            \\
%     \verb|{\'u}|     & {\'u}           \\
%     \verb|{\aa}|     & {\aa}           \\\hline
%   \end{tabular}
%   \begin{tabular}{lc}
%     \hline
%     \textbf{Command} & \textbf{Output} \\
%     \hline
%     \verb|{\c c}|    & {\c c}          \\
%     \verb|{\u g}|    & {\u g}          \\
%     \verb|{\l}|      & {\l}            \\
%     \verb|{\~n}|     & {\~n}           \\
%     \verb|{\H o}|    & {\H o}          \\
%     \verb|{\v r}|    & {\v r}          \\
%     \verb|{\ss}|     & {\ss}           \\
%     \hline
%   \end{tabular}
%   \caption{Example commands for accented characters, to be used in, \emph{e.g.}, Bib\TeX{} entries.}
%   \label{tab:accents}
% \end{table}

% As much as possible, fonts in figures should conform
% to the document fonts. See Figure~\ref{fig:experiments} for an example of a figure and its caption.

% Using the \verb|graphicx| package graphics files can be included within figure
% environment at an appropriate point within the text.
% The \verb|graphicx| package supports various optional arguments to control the
% appearance of the figure.
% You must include it explicitly in the \LaTeX{} preamble (after the
% \verb|\documentclass| declaration and before \verb|\begin{document}|) using
% \verb|\usepackage{graphicx}|.

% \begin{figure}[t]
%   \includegraphics[width=\columnwidth]{example-image-golden}
%   \caption{A figure with a caption that runs for more than one line.
%     Example image is usually available through the \texttt{mwe} package
%     without even mentioning it in the preamble.}
%   \label{fig:experiments}
% \end{figure}

% \begin{figure*}[t]
%   \includegraphics[width=0.48\linewidth]{example-image-a} \hfill
%   \includegraphics[width=0.48\linewidth]{example-image-b}
%   \caption {A minimal working example to demonstrate how to place
%     two images side-by-side.}
% \end{figure*}

% \subsection{Hyperlinks}

% Users of older versions of \LaTeX{} may encounter the following error during compilation:
% \begin{quote}
% \verb|\pdfendlink| ended up in different nesting level than \verb|\pdfstartlink|.
% \end{quote}
% This happens when pdf\LaTeX{} is used and a citation splits across a page boundary. The best way to fix this is to upgrade \LaTeX{} to 2018-12-01 or later.

% \subsection{Citations}

% \begin{table*}
%   \centering
%   \begin{tabular}{lll}
%     \hline
%     \textbf{Output}           & \textbf{natbib command} & \textbf{ACL only command} \\
%     \hline
%     \citep{Gusfield:97}       & \verb|\citep|           &                           \\
%     \citealp{Gusfield:97}     & \verb|\citealp|         &                           \\
%     \citet{Gusfield:97}       & \verb|\citet|           &                           \\
%     \citeyearpar{Gusfield:97} & \verb|\citeyearpar|     &                           \\
%     \citeposs{Gusfield:97}    &                         & \verb|\citeposs|          \\
%     \hline
%   \end{tabular}
%   \caption{\label{citation-guide}
%     Citation commands supported by the style file.
%     The style is based on the natbib package and supports all natbib citation commands.
%     It also supports commands defined in previous ACL style files for compatibility.
%   }
% \end{table*}

% Table~\ref{citation-guide} shows the syntax supported by the style files.
% We encourage you to use the natbib styles.
% You can use the command \verb|\citet| (cite in text) to get ``author (year)'' citations, like this citation to a paper by \citet{Gusfield:97}.
% You can use the command \verb|\citep| (cite in parentheses) to get ``(author, year)'' citations \citep{Gusfield:97}.
% You can use the command \verb|\citealp| (alternative cite without parentheses) to get ``author, year'' citations, which is useful for using citations within parentheses (e.g. \citealp{Gusfield:97}).

% A possessive citation can be made with the command \verb|\citeposs|.
% This is not a standard natbib command, so it is generally not compatible
% with other style files.

% \subsection{References}

% \nocite{Ando2005,andrew2007scalable,rasooli-tetrault-2015}

% The \LaTeX{} and Bib\TeX{} style files provided roughly follow the American Psychological Association format.
% If your own bib file is named \texttt{custom.bib}, then placing the following before any appendices in your \LaTeX{} file will generate the references section for you:
% \begin{quote}
% \begin{verbatim}
% \bibliography{custom}
% \end{verbatim}
% \end{quote}

% You can obtain the complete ACL Anthology as a Bib\TeX{} file from \url{https://aclweb.org/anthology/anthology.bib.gz}.
% To include both the Anthology and your own .bib file, use the following instead of the above.
% \begin{quote}
% \begin{verbatim}
% \bibliography{anthology,custom}
% \end{verbatim}
% \end{quote}

% Please see Section~\ref{sec:bibtex} for information on preparing Bib\TeX{} files.

% \subsection{Equations}

% An example equation is shown below:
% \begin{equation}
%   \label{eq:example}
%   A = \pi r^2
% \end{equation}

% Labels for equation numbers, sections, subsections, figures and tables
% are all defined with the \verb|\label{label}| command and cross references
% to them are made with the \verb|\ref{label}| command.

% This an example cross-reference to Equation~\ref{eq:example}.

% \subsection{Appendices}

% Use \verb|\appendix| before any appendix section to switch the section numbering over to letters. See Appendix~\ref{sec:appendix} for an example.

% \section{Bib\TeX{} Files}
% \label{sec:bibtex}

% Unicode cannot be used in Bib\TeX{} entries, and some ways of typing special characters can disrupt Bib\TeX's alphabetization. The recommended way of typing special characters is shown in Table~\ref{tab:accents}.

% Please ensure that Bib\TeX{} records contain DOIs or URLs when possible, and for all the ACL materials that you reference.
% Use the \verb|doi| field for DOIs and the \verb|url| field for URLs.
% If a Bib\TeX{} entry has a URL or DOI field, the paper title in the references section will appear as a hyperlink to the paper, using the hyperref \LaTeX{} package.

% \section*{Limitations}

% This document does not cover the content requirements for ACL or any
% other specific venue.  Check the author instructions for
% information on
% maximum page lengths, the required ``Limitations'' section,
% and so on.

% \section*{Acknowledgments}

% This document has been adapted
% by Steven Bethard, Ryan Cotterell and Rui Yan
% from the instructions for earlier ACL and NAACL proceedings, including those for
% ACL 2019 by Douwe Kiela and Ivan Vuli\'{c},
% NAACL 2019 by Stephanie Lukin and Alla Roskovskaya,
% ACL 2018 by Shay Cohen, Kevin Gimpel, and Wei Lu,
% NAACL 2018 by Margaret Mitchell and Stephanie Lukin,
% Bib\TeX{} suggestions for (NA)ACL 2017/2018 from Jason Eisner,
% ACL 2017 by Dan Gildea and Min-Yen Kan,
% NAACL 2017 by Margaret Mitchell,
% ACL 2012 by Maggie Li and Michael White,
% ACL 2010 by Jing-Shin Chang and Philipp Koehn,
% ACL 2008 by Johanna D. Moore, Simone Teufel, James Allan, and Sadaoki Furui,
% ACL 2005 by Hwee Tou Ng and Kemal Oflazer,
% ACL 2002 by Eugene Charniak and Dekang Lin,
% and earlier ACL and EACL formats written by several people, including
% John Chen, Henry S. Thompson and Donald Walker.
% Additional elements were taken from the formatting instructions of the \emph{International Joint Conference on Artificial Intelligence} and the \emph{Conference on Computer Vision and Pattern Recognition}.

% Bibliography entries for the entire Anthology, followed by custom entries
%\bibliography{custom,anthology-overleaf-1,anthology-overleaf-2}

% Custom bibliography entries only
\bibliography{custom}

% \appendix

% \section{Example Appendix}
% \label{sec:appendix}

% This is an appendix.

\end{document}